\begin{table}[t]
\centering
\small
\caption{Selective-prediction metrics at fixed coverage in the final scaling audit. Lower risk, AURC, and E-AURC are better; higher failure capture is better.}
\label{tab:selective-fixed}
\begin{threeparttable}
\begin{adjustbox}{width=\textwidth}
\begin{tabular}{llcccccc}
\toprule
Subset & Signal & Risk@80\% & Risk@60\% & Risk@40\% & AURC & E-AURC & Failure capture@20\% abstain \\
\midrule
All seven & Confidence/option & 0.264 & 0.208 & 0.157 & 0.191 & 0.099 & 0.423 \\
All seven & Conf./option + PCA hidden & 0.262 & 0.204 & 0.155 & 0.188 & 0.096 & 0.427 \\
Competent subset & Confidence/option & 0.190 & 0.129 & 0.074 & 0.115 & 0.070 & 0.456 \\
Competent subset & Conf./option + PCA hidden & 0.188 & 0.124 & 0.071 & 0.110 & 0.065 & 0.461 \\
\bottomrule
\end{tabular}
\end{adjustbox}

\begin{tablenotes}[flushleft]\footnotesize\item \parbox{0.95\textwidth}{Metrics are averaged over model-dataset conditions, with the best probe within each reported signal family selected by validation AURC in the earlier audit and by family averages in the Gemma/Llama run. The conclusion is unchanged: hidden augmentation slightly improves some operating points but does not change the practical rule.}\end{tablenotes}
\end{threeparttable}
\end{table}
